% 中文摘要
\begin{abstractzh}
线上部署的智能体大语言模型通过周期性强化学习更新持续进化：每隔一段时间，用新采集的交互数据精炼策略。这本质上是\textbf{持续强化学习}问题——模型必须在习得新能力的同时保住已有能力。与实验室一次性集齐全分布数据不同，线上数据更新快、到达不均衡，往往只能\textbf{按能力域分批、顺序地}训练。这恰好放大了\textbf{灾难性遗忘}：后训练的域会系统性侵蚀先训练的域已掌握的能力。

本文围绕经验回放这一主流抗遗忘手段，面向多能力域的智能体持续强化学习场景，设计并实现一套训练与评测系统。核心贡献有三：

其一，\textbf{按能力分区的九桶经验回放缓冲区}。桶按能力域划分（而非难度），由评测基准官方类目合并为九个纯文本能力域；每桶设固定容量上限与保护下限，容量按任务数平方根加权分配；桶内淘汰、禁止跨桶挤出；桶间按能力空间\textbf{距离}采样回放——离当前训练批越远的桶遗忘风险越高、回放权重越大。桶内采样随机，淘汰采用先进先出。

其二，\textbf{面向持续学习的复合训练目标}，以零框架改动注入GRPO：策略梯度之上叠加逆向KL锚定、经验回放损失、固定开启的熵正则。奖励由单一外部冻结大模型裁判统一打分，采用四维乘性聚合、安全为二元门控。

其三，\textbf{按桶分组的抗遗忘评测协议}：全部能力域训完后一次统一评测，评测集按同样桶分组、组间序与训练序对齐，逐桶给分；再后置地施加随训练先后衰减的权重，量化"越早训练越遗忘"这一持续学习核心指标，并与无持续学习机制的基线对照。

本文在Qwen3.5-9B智能体上跨九个能力域实现并验证整套系统。
\end{abstractzh}
